\unit{Exposición 212 -- Preesquemas cociente}
\sid{212}
\src{99}
\seminares{13.er año, 1960/61, n.º 212\\Febrero de 1961}
\paperhead{TÉCNICAS DE CONSTRUCCIÓN Y TEOREMAS DE EXISTENCIA EN GEOMETRÍA ALGEBRAICA\\
III. PREESQUEMAS COCIENTE}%
  {por Alexander GROTHENDIECK}

\subsection*{Introducción}
\addcontentsline{toc}{subsection}{Introducción}

Los problemas tratados en la presente exposición difieren de los considerados en
los dos anteriores en que se intenta representar ciertos funtores covariantes,
y no ya contravariantes, de esquemas variables. Sin embargo, el procedimiento
de paso al cociente es esencial en muchos problemas de construcción en
geometría algebraica, incluidos los de las Exposiciones I y II ([1], [2]). Así, la
cuestión de la \emph{efectividad de un dato de descenso} sobre un
$T$-preesquema $X$, con respecto a un morfismo fielmente plano y cuasicompacto
$T\longrightarrow S$, equivale a la cuestión de la existencia de un cociente
de $X$ (que satisfaga las propiedades razonables examinadas más adelante) por
la relación de equivalencia plana en $X$ definida por el dato de descenso; las
cuestiones planteadas en [1], A, \S~2 c se resolverán sin duda al mismo tiempo
que las cuestiones formuladas en el párrafo 2 de la presente exposición. Del mismo
modo, el \emph{esquema de Picard} de un $S$-esquema $X$ (para la definición,
véase [2], C, \S~3) puede definirse de diversas maneras como cociente de otros
esquemas (de divisores positivos, o de inmersiones en un esquema proyectivo)
por relaciones de equivalencia planas, mientras que la definición y la
construcción de esos esquemas auxiliares son técnicamente más sencillas: son
esquemas de la forma $\underline{\operatorname{Hom}}_S(X,Y)$ y variantes
definidas en [2], C, \S~2, cuya construcción será objeto de la exposición siguiente
(bajo hipótesis adecuadas de proyectividad). Combinando, pues, los resultados
de la presente exposición y de la siguiente se obtiene la construcción de los esquemas
de Picard bajo hipótesis adecuadas.

El problema del paso al cociente en la categoría de los preesquemas presenta
aún varias cuestiones no resueltas. La más importante se menciona en el
párrafo 8. Actualmente es el único obstáculo para la construcción de
\emph{esquemas de módulos sobre los enteros para curvas de género arbitrario,
variedades abelianas polarizadas}, etc. Esto muestra que su solución merece
los esfuerzos de los especialistas en grupos algebraicos.

\src{100}
\subsection*{1. Relaciones de equivalencia y relaciones de equivalencia efectivas}
\addcontentsline{toc}{subsection}{1. Relaciones de equivalencia y relaciones de equivalencia efectivas}

Sean $\mathcal C$ una categoría y $X$ un objeto de $\mathcal C$. Un par de
morfismos
\[
  p_1,p_2:R\rightrightarrows X
\]
se llama \emph{par de equivalencia} en $\mathcal C$, con \emph{destino} $X$
y \emph{fuente} $R$, si, para todo objeto $T$ de $\mathcal C$, las dos
aplicaciones correspondientes
\[
  p_1(T),p_2(T):R(T)\rightrightarrows X(T)
\]
(donde, para todo objeto $Y$ de $\mathcal C$, se pone
$Y(T)=\operatorname{Hom}(T,Y)$) definen una aplicación
\[
  R(T)\longrightarrow X(T)\times X(T)
\]
que induce una biyección de $R(T)$ sobre el grafo de una relación de
equivalencia en el conjunto $X(T)$. Existe una relación de equivalencia
evidente entre los pares de equivalencia con destino $X$; una clase de
equivalencia para esta última relación se llama una \emph{relación de
equivalencia de $\mathcal C$ en $X$}, o simplemente una relación de
equivalencia cuando no pueda producirse confusión.

Si existe $X\times X$, dar una relación de equivalencia en $X$ equivale a dar
un subobjeto $R$ de $X\times X$ tal que, para todo objeto $T$ de $\mathcal C$,
el subconjunto de
\[
  (X\times X)(T)=X(T)\times X(T)
\]
correspondiente a $R(T)$ sea el grafo de una relación de equivalencia en
$X(T)$. Denotando por $p_1$ y $p_2$ los morfismos de $R$ en $X$ inducidos por
las proyecciones $\operatorname{pr}_1$ y $\operatorname{pr}_2$, la condición
anterior afirma que $(p_1,p_2)$ es un par de equivalencia. También pueden
expresarse diagramáticamente en $\mathcal C$ (siempre que existan $X\times X$
y el producto fibrado $(R,p_2)\times_X(R,p_1)$) los axiomas conjuntistas de
una relación de equivalencia para los $R(T)$ en los $X(T)$, conforme al
principio general explicado en [2], A, \S~1. No lo necesitaremos.

Siempre que se tenga un par de morfismos $(p_1,p_2)$ con la misma fuente $R$
y el mismo destino $X$, puede definirse el \emph{conúcleo} del par como el
objeto $Y$ de $\mathcal C$ que representa el funtor covariante en $Z$
\[
  \operatorname{Hom}_{p_1,p_2}(X,Z),
\]
el conjunto de morfismos $u$ de $X$ en $Z$ tales que $up_1=up_2$. Si existe
$Y$, queda determinado salvo isomorfismo único. Lo denotamos por
$X/(p_1,p_2)$ o, por abuso de notación, por $X/R$; esta última notación se usa
principalmente cuando
\src{101}$(p_1,p_2)$ es un par de equivalencia: entonces es habitual
identificar en la notación la relación de equivalencia definida por el par con
$R$. Obsérvese que, si $Y$ se considera como cociente de $X$, en realidad solo
depende de la relación de equivalencia definida por el par $(p_1,p_2)$.

Partamos ahora de un morfismo
\[
  f:X\longrightarrow Y,
\]
que permite considerar $X$ como un «objeto sobre $Y$», y supongamos que existe
el producto fibrado
\[
  R(f)=X\times_Y X
\]
Sean $p_1$ y $p_2$ sus dos proyecciones. Entonces $(p_1,p_2)$ es un par de
equivalencia, llamado par \emph{asociado} al morfismo $f$. Define, por tanto,
una relación de equivalencia, llamada relación asociada a $f$.

Un par de morfismos $(p_1,p_2)$ con destino $X$ y fuente común $R$ se llama
\emph{par de equivalencia efectivo} si
\begin{enumerate}
  \item[(i)] existe el conúcleo $X/(p_1,p_2)=Y$;
  \item[(ii)] existe el producto fibrado $X\times_Y X$;
  \item[(iii)] el morfismo $R\longrightarrow X\times_Y X$ cuyas componentes son
  $p_1$ y $p_2$ es un isomorfismo.
\end{enumerate}
Entonces $(p_1,p_2)$ es efectivamente un par de equivalencia. La relación de
equivalencia que define se llama también \emph{relación de equivalencia
efectiva}.

Un morfismo $f:X\longrightarrow Y$ se llama \emph{epimorfismo efectivo} si
\begin{enumerate}
  \item[(i)] existe el producto fibrado $X\times_Y X=R$;
  \item[(ii)] existe el cociente $X/(p_1,p_2)$, donde $p_1$ y $p_2$ son las dos
  proyecciones de $R$ en $X$;
  \item[(iii)] el morfismo $X/(p_1,p_2)\longrightarrow Y$ inducido por $f$ es un
  isomorfismo.
\end{enumerate}
Entonces $f$ es efectivamente un epimorfismo, e incluso un epimorfismo
estricto (véase [1], A, \S~2.3); el recíproco vale si existe el producto
fibrado $X\times_Y X$. El objeto cociente de $X$ definido por el epimorfismo
$f$ se llama también \emph{cociente efectivo} de $X$.

Las definiciones anteriores implican la siguiente «correspondencia de
Galois»:

\src{102}
\result{PROPOSICIÓN}{1.1} Existe una correspondencia biunívoca que conserva
los órdenes naturales entre el conjunto de relaciones de equivalencia
efectivas $R$ en $X$ y el conjunto de cocientes efectivos $Y$ de $X$: a $R$
le corresponde el cociente efectivo $X/R$, y a $Y$ le corresponde la relación
de equivalencia efectiva definida por la proyección canónica
$X\longrightarrow Y$ (que está definida por el producto fibrado
$X\times_Y X$ provisto de sus dos proyecciones).

En categorías muy buenas (conjuntos, haces de conjuntos, etc.), todo cociente
es efectivo y toda relación de equivalencia es efectiva. Esto deja de ser
cierto en categorías como la categoría de preesquemas sobre un preesquema fijo
$S$, incluso cuando $S$ es el espectro de un cuerpo, y aun cuando uno se
restringe a esquemas finitos sobre $S$. Las cuestiones de efectividad, e
incluso (para los preesquemas no finitos sobre $S$) las cuestiones de
existencia de cocientes, resultan casi siempre delicadas.

\subsection*{2. Ejemplo: preesquemas finitos sobre $S$}
\addcontentsline{toc}{subsection}{2. Ejemplo: preesquemas finitos sobre S}

Sea $\mathcal C$ la categoría de los preesquemas finitos sobre $S$, donde se
supone que $S$ es localmente noetheriano. Es equivalente a la opuesta de la
categoría de haces coherentes de álgebras conmutativas sobre $S$, o, cuando
$S$ es afín de anillo $A$, a la opuesta de la categoría de $A$-álgebras
finitas sobre $A$ (es decir, finitamente generadas como $A$-módulos). Se sigue
inmediatamente que existen en $\mathcal C$ los límites proyectivos finitos y
los límites inductivos finitos. Esto es bien conocido (sin hipótesis alguna
de finitud\ldots) para los primeros. Así, el producto fibrado de preesquemas
$X,Y$ sobre $S$ corresponde al producto tensorial $B\otimes_A C$ de las
álgebras correspondientes; el núcleo de dos morfismos
$X\rightrightarrows Y$, definidos por dos homomorfismos de $A$-álgebras
$u,v:C\rightrightarrows B$, corresponde al cociente de $B$ por el ideal
generado por los elementos $u(c)-v(c)$, etc. Para los límites inductivos
finitos basta considerar, por una parte, las sumas finitas, que corresponden
al producto ordinario de $A$-álgebras, y, por otra, los conúcleos de pares de
morfismos $X\rightrightarrows Y$, que corresponden (como se ve de inmediato)
al subanillo de $C$ formado por los elementos sobre los que coinciden los
homomorfismos $u,v:C\rightrightarrows B$ (este subanillo es finito sobre $A$
por la hipótesis noetheriana). Además, usando la hipótesis noetheriana, puede
demostrarse que los límites inductivos finitos, y en particular los cocientes,
así construidos en la categoría $\mathcal C$ de preesquemas finitos sobre $S$
son en realidad cocientes en la categoría de \emph{todos} los preesquemas.

\src{103}
Como dijimos en [1], existen epimorfismos no efectivos en $\mathcal C$ (o
epimorfismos no estrictos\footnote{El texto impreso omite «no» antes de
«estrictos». En la categoría considerada existen productos fibrados, y el
párrafo anterior establece precisamente la equivalencia entre efectividad y
epimorfía estricta.}; es lo mismo, puesto que existen los productos fibrados).
\emph{No sé si las relaciones de equivalencia son siempre efectivas} cuando
no se impone ninguna hipótesis de planitud. En esta dirección solo he obtenido
resultados positivos muy parciales, indispensables para la demostración del
teorema fundamental de la teoría formal de módulos (véase [2], B, teo. 1).
Señalemos que en el problema planteado es fácil reducirse al caso en que $S$
es el espectro de un anillo local artiniano $A$ con cuerpo residual
algebraicamente cerrado. Pero incluso si $A$ es un cuerpo, no se conoce la
respuesta.

También puede considerarse un preesquema $X$ sobre $S$ que ya no se suponga
finito sobre $S$, pero provisto de una relación de equivalencia $R$ en $X$ tal
que $p_1:R\longrightarrow X$ sea un morfismo finito. Entonces $R$ se llama
una relación de equivalencia finita. Supongamos, para simplificar, que $S$ y
$X$ son afines (y, por tanto, $R$ es afín, de modo que la situación se reduce
a una cuestión de álgebra conmutativa pura). \emph{Incluso en este caso se
ignora si existe un cociente $X/R=Y$ y si el morfismo canónico
$X\longrightarrow Y$ es finito}. (El caso más sencillo es aquel en que $S$ es
el espectro de un cuerpo $k$ y $X$ es el espectro de $k[t]$, es decir, la
recta afín.) Por supuesto, si los dos problemas anteriores tienen soluciones
afirmativas, puede concluirse en la situación presente que $R$ es efectiva.
Obsérvese que el problema de la \emph{existencia} de un cociente $Y$ y de la
finitud de $f:X\longrightarrow Y$ se plantea exactamente en los mismos
términos si, en lugar del grafo de una relación de equivalencia en $X$, solo
se tiene un pregrafo de equivalencia en $X$ en el sentido del párrafo 4.

La cuestión del paso al cociente por una relación de equivalencia finita más
o menos arbitraria se presenta al construir preesquemas «pegando» preesquemas
dados $X_i$ a lo largo de ciertos subpreesquemas cerrados; la ley de pegado se
expresa precisamente mediante una relación de equivalencia finita en el
preesquema suma $X$ de los $X_i$. Cabe esperar también que la solución de los
problemas aquí planteados, y de diversas variantes, sea una condición previa
para desarrollar una técnica general de construcciones no proyectivas en la
dirección iniciada en [2].

El único hecho positivo general conocido por el autor es el siguiente:

\result{PROPOSICIÓN}{2.1} Sea $S$ un preesquema localmente noetheriano, sea
$s$ un punto de $S$ y sea $\Omega$ una extensión algebraicamente cerrada de
$k(s)$. Considérese el
\src{104}«funtor fibra» correspondiente $F$, que asigna a todo $S$-esquema
$X$ finito sobre $S$ el conjunto de puntos con valores en $\Omega$ de $X_s$.
Este funtor (trivialmente exacto por la izquierda) es \emph{exacto por la
derecha}, es decir, conmuta con los límites inductivos finitos y, en
particular, con los conúcleos de pares de morfismos.

Aplicando este resultado a todos los «puntos geométricos» de $S$, se deduce
que la categoría «cociente» $\mathcal C'$ de $\mathcal C$, obtenida razonando
«módulo morfismos radiciales sobreyectivos» (es decir, adjuntando formalmente
inversos a los morfismos en cuestión), es una categoría «geométrica», esto es,
satisface las mismas propiedades «de naturaleza finita» que la categoría de
conjuntos. En particular, toda relación de equivalencia en ella es efectiva.
Esto implica que, si $R$ es una relación de equivalencia en un $X$ finito
sobre $S$, entonces el morfismo canónico
\[
  R\longrightarrow X\times_Y X\qquad (\text{donde }Y=X/R)
\]
es \emph{radicial y sobreyectivo} (de hecho, una inmersión cerrada
sobreyectiva, puesto que es un monomorfismo).

\subsection*{3. El caso de un grupo de operadores}
\addcontentsline{toc}{subsection}{3. El caso de un grupo de operadores}

Supongamos de nuevo que $\mathcal C$ es una categoría arbitraria. Sean $G,X$
objetos de $\mathcal C$, y supongamos que $G$ es un grupo de operadores de
$\mathcal C$ sobre el objeto $X$ de $\mathcal C$. Esto significa (véase [2],
A, \S~1) que, para todo objeto $T$ de $\mathcal C$, se han dado una estructura
de grupo en $G(T)$ y una estructura de $G(T)$-conjunto en $X(T)$, de tal modo
que, al variar $T$, las estructuras en cuestión «varían funtorialmente» con
$T$. Si existen los productos $G\times G$ y $G\times X$ en $\mathcal C$, una
estructura tal puede definirse también mediante un par de morfismos
\[
  G\times G\longrightarrow G,
  \qquad \pi:G\times X\longrightarrow X,
\]
sometidos a la condición de que, para todo objeto $T$ de $\mathcal C$, las
leyes de composición correspondientes en los conjuntos $G(T)$ y $X(T)$ hagan
de $G(T)$ un grupo que actúa en $X(T)$. Traducir este axioma en la
conmutatividad de ciertos diagramas en $\mathcal C$ es fácil pero tedioso y,
de hecho, completamente inútil en todos los casos que conozco.

Supongamos que existe $G\times X$ y consideremos los dos morfismos
\[
  p_1,p_2:G\times X\rightrightarrows X
\]
dados por\footnote{El texto impreso da \texttt{pr\_1}. Puesto que los dos
morfismos tienen codominio $X$, el primero de ellos debe ser
\texttt{pr\_2}; se corrige aquí la errata tipográfica.}
\src{105}
\[
  p_1=\operatorname{pr}_2,\qquad p_2=\pi.
\]
Se comprueba de inmediato que $(p_1,p_2)$ es un par de equivalencia si y solo
si, para todo objeto $T$ de $\mathcal C$, la aplicación
\[
  G(T)\times X(T)\simeq(G\times X)(T)
  \longrightarrow X(T)\times X(T),
\]
definida por este par es inyectiva, es decir, si el grupo $G(T)$ actúa
\emph{libremente} en el conjunto $X(T)$, esto es, si $g\in G(T)$,
$x\in X(T)$ y $g\cdot x=x$ implican que $g$ es el elemento identidad de
$G(T)$. Se dice entonces que $G$ \emph{actúa libremente en $X$} (o que $X$ es
un \emph{espacio principal de $\mathcal C$ bajo $G$}). La relación de
equivalencia asociada al par $(p_1,p_2)$ se llama entonces \emph{relación de
equivalencia definida por el grupo $G$ que actúa libremente en $X$}. Si existe
también $X\times X$ y se considera el morfismo
\[
  p:G\times X\longrightarrow X\times X
\]
definido por el par $(p_1,p_2)$, la condición de que $G$ actúe libremente
significa que $p$ es un \emph{monomorfismo}.

Por supuesto, incluso cuando $G$ no actúa libremente en $X$, se desean
criterios para la existencia de un cociente de $X$ por $G$, es decir, del
conúcleo del par anterior $(p_1,p_2)$.

El conúcleo en cuestión se denotará a menudo por $X/G$, o preferiblemente por
$G\backslash X$ si $G$ actúa por la izquierda (la primera notación se reserva
para el caso en que $G$ actúa por la derecha). Obsérvese que, incluso cuando
existe la «imagen» de $G\times X$ por $p$ (definida, por ejemplo, como el
menor subobjeto de $X\times X$ a través del cual se factoriza $p$), digamos
$R$, esta no es casi nunca una relación de equivalencia en $X$. Si se intenta
entonces pasar directamente al cociente por $R$ (o, más precisamente, por el
par de morfismos de $R$ en $X$ inducidos por las dos proyecciones
$\operatorname{pr}_i$), se pierden irremediablemente las características
especiales del par original $(p_1,p_2)$. Por ello es importante hallar una
generalización de la noción de relación de equivalencia que se aplique
directamente al par definido por un grupo de operadores de $\mathcal C$.

\subsection*{4. Prerrelaciones de equivalencia}
\addcontentsline{toc}{subsection}{4. Prerrelaciones de equivalencia}

Recordemos que una categoría en la que todo morfismo es un isomorfismo se
llama un \emph{grupoide}. Por otra parte, una categoría debe definirse como
constituida por dos conjuntos subyacentes $(X,R)$, el conjunto de
\emph{objetos} y el conjunto de \emph{flechas}, provistos de las estructuras
siguientes:
\src{106}
\begin{enumerate}
\item[(i)] Un par de aplicaciones
\[
  p_1,p_2:R\rightrightarrows X
\]
llamadas \emph{aplicación fuente} y \emph{aplicación destino};

\item[(ii)] Una aplicación
\[
  \pi:(R,p_2)\times_X(R,p_1)\longrightarrow R
\]
llamada \emph{aplicación de composición}.
\end{enumerate}
Estos datos deben satisfacer axiomas conocidos, que no repetiremos aquí y que
podrían expresarse mediante la conmutatividad de ciertos diagramas, junto con
la existencia de una aplicación $D:X\to R$ (necesariamente única) que haga
conmutativos otros dos diagramas. La aplicación $D$ envía un objeto a la
flecha identidad correspondiente y satisface
\[
  p_1D=p_2D=\operatorname{id}_X.
\]
Decir que la categoría es un grupoide significa entonces que existe una
aplicación
\[
  s:R\longrightarrow R
\]
(necesariamente única), llamada la \emph{inversión} de $R$, que transforma
cada flecha en una flecha inversa. Esto podría expresarse mediante la
conmutatividad de otros cuatro diagramas formados a partir de $s$, $\Delta$ y
los datos anteriores; los dos primeros son
\[
  p_1s=p_2,\qquad p_2s=p_1.
\]

Recordados estos puntos, las definiciones generales de ([2], A, \S~1)
muestran, en particular, qué se entiende por una estructura de
$\mathcal C$-\emph{categoría}, respectivamente de
$\mathcal C$-\emph{grupoide}, sobre un par de objetos $X,R$ de una categoría
arbitraria $\mathcal C$: por definición, es el dato, para todo objeto $T$ de
$\mathcal C$, de una estructura de categoría, respectivamente de grupoide, en
el sentido conjuntista, cuyo conjunto de objetos es $X(T)$ y cuyo conjunto de
flechas es $R(T)$, «variando funtorialmente» estas estructuras con la variable
$T$. Esto proporciona, por tanto, dos morfismos
\[
  p_1,p_2:R\rightrightarrows X
\]
llamados \emph{morfismo fuente} y \emph{morfismo destino}, y, cuando existe el
producto fibrado requerido, un morfismo
\src{107}
\[
  \pi:(R,p_2)\times_X(R,p_1)\longrightarrow R
\]
llamado \emph{morfismo de composición}. Estos tres datos bastan entonces para
determinar la estructura de categoría sobre $X,R$, y el axioma que se les
impone es el siguiente: para todo $T$, los tres datos correspondientes para
$X(T),R(T)$ definen una estructura de categoría (respectivamente de grupoide)
sobre este par de conjuntos. Cuando proceda, esto puede expresarse mediante
la conmutatividad de ciertos diagramas en los que interviene un morfismo
determinado de manera única
\[
  D:X\longrightarrow R
\]
y, en el caso de los grupoides, un morfismo determinado de manera única
\[
  s:R\longrightarrow R;
\]
los diagramas se escriben exactamente como en el caso «conjuntista». Esta
laboriosa interpretación de los axiomas es, afortunadamente, inútil en la
práctica. El único interés teórico de poder expresar los datos y axiomas en
términos de morfismos e igualdades de morfismos entre ciertos productos
fibrados es el siguiente: si $F:\mathcal C\to\mathcal C'$ es un funtor exacto
por la izquierda (es decir, que conmuta con productos finitos y productos
fibrados), entonces transforma una $\mathcal C$-categoría (respectivamente un
$\mathcal C$-grupoide) en una $\mathcal C'$-categoría (respectivamente un
$\mathcal C'$-grupoide), siempre que existan productos finitos y productos
fibrados en $\mathcal C$.

En la práctica, es importante saber interpretar los morfismos
$p_1,p_2,\pi,D,s$ como \emph{operaciones simpliciales} en un objeto
semisimplicial o simplicial adecuado de $\mathcal C$, al menos cuando existen
productos fibrados en $\mathcal C$. Para fijar la terminología, introduzcamos
la \emph{categoría $\mathcal S$ de los símplices estándar} como la categoría
cuyos objetos son los conjuntos finitos de la forma
\[
  \Delta_n=[0,n],\qquad\text{donde }n\geq 0
\]
(el intervalo de enteros de $0$ a $n$), y cuyos morfismos o flechas son
\emph{aplicaciones arbitrarias} entre esos conjuntos finitos. Obsérvese que la
categoría $\mathcal S$ es equivalente a la categoría de conjuntos finitos
\emph{no vacíos}, con las aplicaciones de conjuntos finitos como morfismos.
En $\mathcal S$ existe claramente el coproducto de una familia finita
\emph{no vacía} de objetos, así como la suma amalgamada de dos objetos bajo un tercero
(la operación dual del producto fibrado). Denotemos por $\mathcal S'$ la
subcategoría de $\mathcal S$ que tiene los mismos objetos pero cuyos morfismos
son las \emph{aplicaciones crecientes} entre los $\Delta_n$. Esta categoría es
equivalente a la categoría de conjuntos finitos no vacíos totalmente
ordenados. En esta categoría el
\src{108}
coproducto de dos objetos nunca existe, y la suma amalgamada de dos objetos $A,B$ bajo
un tercero $C$ no existe en general (tómese, por ejemplo, $C=\Delta_0$,
$A=B=\Delta_1$, con las dos aplicaciones estructurales $u:C\to A$ y
$v:C\to B$ iguales). Sin embargo, existe en ciertos casos; por ejemplo,
\[
  C=\Delta_0,\qquad A=\Delta_m,\qquad B=\Delta_n,\qquad
  u(0)=m,\qquad v(0)=0,
\]
y en este caso
\[
  \Delta_m\amalg_{\Delta_0}\Delta_n=\Delta_{m+n}.
\]

Un \emph{objeto simplicial} (respectivamente, un \emph{objeto
semisimplicial}) en una categoría $\mathcal C$ es, por definición, un funtor
contravariante $K$ de $\mathcal S$ (respectivamente de $\mathcal S'$) en
$\mathcal C$. Por tanto, un objeto simplicial define por restricción un objeto
semisimplicial, pero la primera noción difiere esencialmente de la segunda por
la presencia de \emph{operaciones de simetría} en los $K_n=K(\Delta_n)$,
correspondientes a las imágenes por el funtor $K$ de los elementos del grupo
simétrico de $n+1$ letras (considerado como el grupo de automorfismos de
$\Delta_n$ en $\mathcal S$).

Para cada $n$, sea $\Delta'_n$ (respectivamente $\Delta''_n$) la categoría
finita cuyo conjunto de objetos es $\Delta_n$, definida por la relación de
orden caótica (respectivamente por el orden total natural) en $\Delta_n$ (el
conjunto de flechas es el grafo de esa relación de orden). Está claro que
$\Delta'_n$ (respectivamente $\Delta''_n$) depende funtorialmente del objeto
$\Delta_n$ de $\mathcal S$ (respectivamente de $\mathcal S'$). Si $Z$ es una
categoría, entonces $\operatorname{Hom}(\Delta'_n,Z)$ (respectivamente
$\operatorname{Hom}(\Delta''_n,Z)$), al variar $\Delta_n$, es un funtor de
$\mathcal S$ (respectivamente de $\mathcal S'$) a la categoría de conjuntos,
es decir, un conjunto simplicial (respectivamente un conjunto
semisimplicial), que se dice \emph{asociado a la categoría $Z$}; denotémoslos
por $Z'$ y $Z''$. Hay, además, un homomorfismo natural evidente del conjunto
semisimplicial asociado a $Z'$ hacia $Z''$, y este último es un isomorfismo si
y solo si $Z$ es un grupoide. Entendido esto:

\result{PROPOSICIÓN}{4.1} El funtor $Z\longmapsto Z''$ de la categoría de
categorías a la categoría de conjuntos semisimpliciales es plenamente fiel y
define una equivalencia entre la categoría de las \emph{categorías} y la
categoría de los conjuntos semisimpliciales, es decir, los funtores
contravariantes $K$ de $\mathcal S'$ en $(\mathrm{Set})$ que transforman las
sumas amalgamadas $A\amalg_C B$ del tipo especificado anteriormente en productos
fibrados de conjuntos. Asimismo, el funtor $Z\longmapsto Z'$ de la categoría
de grupoides a la categoría de conjuntos simpliciales es plenamente fiel y
define una equivalencia entre la categoría de los \emph{grupoides} y la
categoría de los
\src{109}
conjuntos simpliciales, es decir, los funtores contravariantes $K$ de
$\mathcal S$ en $(\mathrm{Set})$ que transforman las sumas amalgamadas en productos
fibrados.

Así, las categorías pueden considerarse como conjuntos semisimpliciales
particulares, y los grupoides como conjuntos simpliciales particulares,
siempre que, por supuesto, se razone «salvo isomorfismo», como exige la
interpretación de ciertas estructuras en términos de otras. El procedimiento
usual de reducción al caso conjuntista implica entonces:

\result{COROLARIO}{4.2} El enunciado anterior sigue siendo válido cuando las
categorías, los grupoides y los conjuntos simpliciales se sustituyen por
$\mathcal C$-categorías, $\mathcal C$-grupoides y objetos simpliciales en
$\mathcal C$, siempre que existan productos fibrados en $\mathcal C$.

El objeto semisimplicial $K$ en $\mathcal C$ asociado a una categoría
$(X,R,\ldots)$ en $\mathcal C$ puede hacerse explícito interpretando la
componente $K_n=K(\Delta_n)$ de $K$ como la potencia fibrada $n$-ésima de
$(R,p_1)$ sobre $X$, o, mejor, mediante la fórmula de recurrencia
\[
  K_0=X,\qquad
  K_n=(K_{n-1},p_{n-1}^{(n-1)})\times_X(R,p_1),
\]
donde los $p_i^{(n-1)}$ ($0\leq i\leq n-1$) son las proyecciones naturales de
$K_{n-1}$ en $X$ (que se definen también recursivamente).\footnote{El texto impreso lleva $(n+1)$, $K_0=R$ y $p_n^{(n-1)}$ en estos tres lugares. La indexación del nervio exige $n$, $K_0=X$ y $p_{n-1}^{(n-1)}$, respectivamente.} De este modo,
$p_1,p_2,\pi,D,s$ se interpretan como las operaciones simpliciales
correspondientes respectivamente a los morfismos siguientes de $\mathcal S$:
la cara $0$ de $\Delta_1$, la cara $1$ de $\Delta_1$, la cara $(0,2)$ de
$\Delta_2$, la degeneración $\Delta_1\to\Delta_0$ y la simetría de
$\Delta_1$. Todas las demás operaciones semisimpliciales (respectivamente,
simpliciales) se obtienen formalmente a partir de las cuatro (respectivamente,
cinco) anteriores mediante composición y productos fibrados.

Una \emph{prerrelación de equivalencia} en un objeto $X$ de una categoría se
define ahora como el dato de un grupoide cuyo objeto de objetos (por así
decirlo) es $X$. Un dato tal comprende, por tanto, entre otras cosas, un
objeto $R$ y dos morfismos
\[
  p_1,p_2:R\rightrightarrows X.
\]
Pero obsérvese que estos datos por sí solos no determinan la estructura en
cuestión, a diferencia de lo que ocurre con los pares de equivalencia. En esta
exposición nos interesa esta noción para obtener criterios de posibilidad del
paso al cociente, es decir, para formar el conúcleo del par $(p_1,p_2)$. Por
ello, el enunciado del problema no hace intervenir los datos adicionales
inherentes a un grupoide. En las demostraciones de los resultados que siguen,
\src{110}
usaremos, no obstante, esos datos adicionales y, en particular, las
operaciones simpliciales (incluidas las operaciones de simetría) hasta la
dimensión $3$ (en las que interviene la cuarta potencia fibrada de $R$ sobre
$X$).

Una relación de equivalencia en un objeto $X$ de $\mathcal C$ define una
prerrelación de equivalencia: basta verlo en el caso conjuntista, donde a una
relación de equivalencia en un conjunto $X$ se asocia el grupoide cuyo
conjunto de objetos es $X$ y cuyo conjunto de flechas es el grafo de la
relación de equivalencia.

Un monoide de $\mathcal C$ $G$ que actúa sobre un objeto $X$ de $\mathcal C$
define una $\mathcal C$-categoría cuyos objetos subyacentes son $R=G\times X$
y $X$ (siempre que exista $G\times X$), y esta es un $\mathcal C$-grupoide si
y solo si $G$ es un grupo. De nuevo, basta comprobarlo en el caso conjuntista.
La composición de las flechas $(g,a)$ y $(g',g\cdot a)$ se define entonces por
\[
  (g',g\cdot a)(g,a)=(g'g,a),
\]
es decir, si $a,b\in X$, entonces $\operatorname{Hom}(a,b)$ es por definición
el transportador de $a$ a $b$, y los morfismos se componen mediante la
composición de los elementos de $G$.

\medskip
\noindent\textbf{OBSERVACIÓN. ---} Las dificultades lógicas planteadas por un
enunciado como (4.1) pueden evitarse entendiendo que todos los objetos
considerados pertenecen a un «Universo» fijo (que es a su vez un conjunto).

\fsection{5}{Cociente por una relación de equivalencia finita y plana}

\result{TEOREMA}{5.1} Sea $X=\operatorname{Spec}(B)$ un esquema afín y sea
$R$ una prerrelación de equivalencia en $X$ cuya componente $R_1$ sea afín,
digamos $R_1=\operatorname{Spec}(C)$. Supongamos que la primera proyección
$p_1:R_1\to X$ es finita y localmente libre; en otras palabras, supongamos
que el homomorfismo de anillos correspondiente $p_1^*:B\to C$ hace de $C$ un
$B$-módulo proyectivo finitamente generado. Sea $A$ el subanillo de $B$ que
es el núcleo del par de homomorfismos
\[
  p_1^*,p_2^*:B\rightrightarrows C
\]
(el conjunto de elementos $b$ tales que $p_1^*(b)=p_2^*(b)$). Sea
$Y=\operatorname{Spec}(A)$ y sea $f:X\to Y$ el morfismo definido por la
inclusión de $A$ en $B$. En estas condiciones:
\begin{enumerate}
\item[(i)] $B$ es entero sobre $A$, es decir, $f$ es un morfismo entero.

\item[(ii)] El morfismo $f$ es sobreyectivo y sus fibras son las clases de
equivalencia
\src{111}
$p_2(p_1^{-1}(x))$ en el conjunto $X$ módulo $R$; además, la topología de $Y$
es el cociente de la de $X$.

\item[(iii)] $Y$ es el cociente de $X$ por $R$ en la categoría de los
preesquemas.

\item[(iv)] Si $R$ proviene de una \emph{relación de equivalencia}, entonces
el morfismo $f:X\to Y$ es finito localmente libre (es decir, $B$ es un
$A$-módulo proyectivo finitamente generado), y la relación de equivalencia es
efectiva, esto es, $R_1\to X\times_Y X$ es un isomorfismo.
\end{enumerate}

Este teorema generaliza el conocido teorema relativo a un grupo finito $G$ que
actúa por automorfismos en el anillo $B$ y al anillo de invariantes $A$, y su
demostración es análoga a la demostración conocida. El enunciado (iii) puede
precisarse como sigue:

\result{COROLARIO}{5.2} El morfismo canónico
$R_1\to X\times_Y X$ es \emph{sobreyectivo}.

Sea todavía $R$ una prerrelación de equivalencia «finita y localmente libre»
en $X$, pero sea ahora $X$ un preesquema arbitrario. Supongamos que puede
hallarse un preesquema $Y$ y un morfismo $f:X\to Y$ tales que $fp_1=fp_2$ y,
además, tales que la sucesión de homomorfismos de haces de anillos sobre $Y$
\[
  \mathcal O_Y\longrightarrow f_*(\mathcal O_X)
  \rightrightarrows g_*(\mathcal O_R)
\]
sea exacta (donde $g=fp_i$). El teorema proporciona entonces conclusiones
análogas a (i)--(iv) anteriores; en particular, por (iii), $Y$ es el cociente
de $X$ por $R$ y, por tanto, queda determinado salvo isomorfismo único. En
estas condiciones se dice que la prerrelación de equivalencia $R$ en $X$ es
\emph{admisible}. Con esta definición:

\result{TEOREMA}{5.3} Sea $X$ un preesquema y sea $R$ una prerrelación de
equivalencia en $X$ tal que $p_1:R_1\to X$ sea finito y localmente libre.
Entonces $R$ es admisible si y solo si toda clase de equivalencia conjuntista
$p_2(p_1^{-1}(x))$ en $X$ módulo $R$ está contenida en un abierto afín
(condición que siempre se satisface si todo subconjunto finito de $X$ está
contenido en un abierto afín, por ejemplo si $X$ es cuasiproyectivo sobre un
esquema afín).

En efecto, se demuestra sin dificultad que toda clase de equivalencia módulo
$R$ en $X$ está entonces contenida en un abierto afín \emph{estable bajo
$R$}, y se construye el cociente $Y$ pegando las piezas obtenidas del Teorema
5.1.
\src{112}

\result{COROLARIO}{5.4} Supongamos que se satisface esta condición y que,
además, $R$ proviene de una relación de equivalencia. Entonces esa relación es
efectiva, es decir, $R_1\to X\times_Y X$ es un isomorfismo, y $f:X\to Y$ es
finito y localmente libre.

Se concluye inmediatamente por «descenso»:

\result{COROLARIO}{5.5} Bajo las hipótesis de (5.4), $X$ tiene rango $n$ en
todo punto sobre $Y$ si y solo si $(R_1,p_1)$ tiene rango $n$ en todo punto
sobre $X$. Si $X,R_1$ son $Z$-preesquemas y $p_1,p_2$ son $Z$-morfismos, de
modo que $Y$ es un $Z$-preesquema, entonces $X$ es plano sobre $Z$ si y solo
si $Y$ es plano sobre $Z$.

En resumen:

\result{ESCOLIO}{} El dato de un morfismo finito, localmente libre y
sobreyectivo $f:X\to Y$ de preesquemas equivale al dato de un preesquema $X$
provisto de una relación de equivalencia $R$ tal que $p_1:R\to X$ sea finito
y localmente libre y toda clase $p_2(p_1^{-1}(x))$ esté contenida en un
abierto afín.

\medskip
\noindent\textbf{OBSERVACIONES 5.6.}

\smallskip
\noindent $1^\circ$ No se ha necesitado ninguna hipótesis noetheriana.

\smallskip
\noindent $2^\circ$ Esta noción de paso al cociente contiene, como caso
particular, el «descenso inseparable» de CARTIER, que consiste en determinar
los morfismos finitos localmente libres $f:X\to Y$ tales que
$f_*(\mathcal O_X)$ admita una $p$-base relativa a $\mathcal O_Y$ (donde $X$
es un preesquema dado cuyo haz $\mathcal O_X$ es anulado por el número primo
$p>0$). Obsérvese que el resultado se expresa fácilmente sin ninguna
hipótesis de regularidad sobre los anillos locales y sin suponer que $X$ sea
un esquema algebraico sobre un cuerpo. La teoría de JACOBSON--BOURBAKI se
obtiene tomando por $X$ el espectro de un cuerpo de característica $p$.

\smallskip
\noindent $3^\circ$ GABRIEL había obtenido anteriormente un caso particular
del Teorema (5.3) en la teoría del paso al cociente para grupos conmutativos
finitos sobre un cuerpo $k$. [Compárese (7.4).]

\fsection{6}{Cociente por una relación de equivalencia propia y plana}

\result{TEOREMA}{6.1} Sea $S$ un preesquema localmente noetheriano, sea $X$
un $S$-esquema cuasiproyectivo y sea $R$ una prerrelación de equivalencia en
$X$ tal que:
\src{113}
\begin{enumerate}
\item[(a)] $p_1:R_1\to X$ es propio y plano;
\item[(b)] $R_1\to X\times_S X$ es finito o, equivalentemente, por (a), tiene
fibras finitas (condición que se satisface automáticamente si $R$ proviene de
una relación de equivalencia).
\end{enumerate}
En estas condiciones:
\begin{enumerate}
\item[(i)] Existe $X/R=Y$ y, si $S$ es noetheriano y $R$ proviene de una
relación de equivalencia, es cuasiproyectivo sobre $S$.\footnote{El texto
impreso no contiene la segunda hipótesis. La fe de erratas vinculada de la Exposición
212 dice que la cuasiproyectividad parece haberse demostrado entonces solo
cuando $R$ proviene de una relación de equivalencia.}

\item[(ii)] El morfismo canónico $f:X\to Y$ es sobreyectivo, propio y abierto;
sus fibras son las clases de equivalencia $p_2(p_1^{-1}(x))$ en $X$ módulo
$R$. Así, $Y$ se identifica con el espacio topológico cociente de $X$ por la
relación de equivalencia conjuntista definida por $R$. Finalmente,
$R_1\to X\times_Y X$ es sobreyectivo.

\item[(iii)] Si $R$ proviene de una relación de equivalencia, entonces esa
relación es efectiva, es decir, $R_1\to X\times_Y X$ es un isomorfismo, y
además $f:X\to Y$ es plano (por tanto, fielmente plano).
\end{enumerate}

Para la demostración, uno se reduce a (5.1) mediante cuasisecciones adecuadas
de $X$ para $R$; la demostración es análoga a la construcción de grupos
algebraicos cociente en el Seminario Chevalley.

En resumen:

\result{ESCOLIO}{} Sea $X$ cuasiproyectivo sobre un $S$ localmente
noetheriano. El dato de un morfismo propio, fielmente plano y sobreyectivo
$f:X\to Y$ de $X$ a un $S$-preesquema $Y$ equivale al dato de una relación de
equivalencia $R$ en $X$ tal que $p_1:R\to X$ sea propio y plano.

El mismo método da el resultado siguiente:

\result{TEOREMA}{6.2} Sea $S$ un preesquema noetheriano, sea $X$ un
preesquema de tipo finito sobre $S$ y sea $R$ una prerrelación de equivalencia
en el $S$-preesquema $X$. Supongamos que:
\begin{enumerate}
\item[(a)] $p_1:R_1\to X$ es plano y de tipo finito;
\item[(b)] el morfismo $R_1\to X\times_S X$ es cuasifinito (es decir, tiene
fibras finitas).
\end{enumerate}
Entonces existe un abierto denso $U$ de $X$ que está \emph{saturado} para $R$
y tal que:
\begin{enumerate}
\item[(i)] Si $R_U$ es la prerrelación de equivalencia inducida por $R$ en
$U$, entonces existe $U/R_U$ y es de tipo finito sobre $S$.

\item[(ii)] El morfismo canónico $U\to U/R_U$ es sobreyectivo y abierto; sus
fibras son las clases de equivalencia conjuntistas para $R_U$ (así,
$U/R_U$ es el espacio topológico cociente de $U$ por la relación de
equivalencia conjuntista definida por $R_U$). Finalmente, el morfismo
\[
  (R_U)_1\longrightarrow U\times_{U/R_U}U
\]
es sobreyectivo.
\src{114}

\item[(iii)] Si $R$ proviene de una relación de equivalencia, puede disponerse
que $U\to U/R_U$ sea fielmente plano y que $R_U$ sea efectiva.
\end{enumerate}

Este es un resultado de naturaleza esencialmente «birracional».

\medskip
\noindent\textbf{OBSERVACIONES 6.3.}

\smallskip
\noindent $1^\circ$ No sé si la hipótesis (b) es innecesaria en (6.1) y
(6.2). Al tomar cocientes por grupos, nos obliga en la práctica a
restringirnos al caso en que todos los estabilizadores son grupos finitos.

\smallskip
\noindent $2^\circ$ Puede preguntarse si existen resultados análogos a (6.1)
y (6.2) sin hipótesis de planitud. No conozco ningún contraejemplo en esta
dirección. Por otra parte, incluso conservando la hipótesis de planitud y
restringiéndose a relaciones de equivalencia para las cuales $p_1:R\to X$ es
plano y cuasifinito (pero no finito), con $X$ afín, puede ocurrir que $R$ no
sea efectiva: tómense las relaciones de equivalencia inducidas en los
subconjuntos abiertos afines que recubren la variedad de Nagata (en la que un
grupo de orden dos actúa de manera «no admisible»).

\fsection{7}{Aplicaciones}

Como dijimos en la introducción, la aplicación más importante de (6.1) es la
construcción de esquemas de Picard, así como la solución de otros diversos
problemas de módulos, sobre los cuales volveremos más adelante.

Se obtiene una demostración sencilla del resultado siguiente de SHIMURA:

\result{PROPOSICIÓN}{7.1} Sea $A$ un esquema abeliano definido sobre un
anillo de valoración discreta $V$ de cuerpo de fracciones $K$. Entonces todo
esquema abeliano $B'$ sobre $K$ que sea isógeno a un cociente de
$A\otimes_V K$ «tiene buena reducción sobre $V$», es decir, es isomorfo a
algún $B\otimes_V K$,\footnote{El texto impreso contiene $B\otimes_V A$.
Como $B$ es un $V$-esquema y la afirmación concierne a su fibra genérica, el
factor de base debe ser $K$.} donde $B$ es un esquema abeliano sobre $V$
(esencialmente único, como recordamos).

Podemos suponer que $B'$ es el cociente de $A_K$ por un subesquema en grupos
$C'$. (N.B.: $C'$ no será en general «reducido», es decir, sus anillos locales
tendrán elementos nilpotentes.) Considérese el subesquema cerrado $C$ de $A$
que es la «clausura» de $C'$, esto es, el menor subesquema cerrado de $A$ tal
que $C_K$ contiene a $C'$. Se tiene $C_K=C'$, y el hecho de que $V$ sea un
anillo de valoración discreta implica fácilmente que $C$ es un subesquema en
grupos de $A$ sobre $V$. Puesto que $A$ es propio sobre
$\operatorname{Spec}(V)=S$, también lo es $C$; además, $A$ es incluso
proyectivo sobre $S$. Podemos, por tanto, aplicar (6.1) para construir
$A/C=B$,
\src{115}
que es el $B$ buscado.

Por último, argumentos esencialmente conocidos permiten deducir de (6.2) el
resultado siguiente:

\result{TEOREMA}{7.2} Sea $S$ el espectro de un anillo artiniano, sean $F$ y
$G$ esquemas en grupos de tipo finito sobre $S$, y sea $u:F\to G$ un
homomorfismo de esquemas en grupos sobre $S$. Supongamos que:
\begin{enumerate}
\item[(i)] $F$ es plano sobre $S$;
\item[(ii)] el núcleo de $u$ es finito.
\end{enumerate}
En estas condiciones existe el esquema cociente $G/F$; el morfismo canónico
$G\to G/F$ es sobreyectivo y abierto, y sus fibras son las clases de
equivalencia conjuntistas definidas por las acciones por la derecha de $F$ en
$G$. Finalmente, si $u$ es un monomorfismo, entonces $G\to G/F$ es plano y
el morfismo
\[
  G\times F\longrightarrow G\times_{G/F}G
\]
es un isomorfismo. En otras palabras, $G$ es un espacio homogéneo principal
sobre $G/F$ con grupo estructural (que actúa por la derecha) $F$, o, más
exactamente, $F\times_S(G/F)$ considerado como esquema en grupos sobre $G/F$
(véase [1], B, \S~6).

\result{COROLARIO}{7.3} En estas condiciones, $G$ es plano sobre $S$ si y
solo si $G/F$ es plano sobre $S$. Si se cumple esta condición, la formación
del cociente por $F$ conmuta con toda extensión de la base $S$; y, si $F$ es
un subgrupo normal de $G$, entonces $G/F$ puede dotarse de la estructura de
\emph{grupo cociente} de $G$ por $F$.

La situación es especialmente sencilla cuando $S$ es el espectro de un
cuerpo, pues entonces todo $S$-preesquema es automáticamente plano sobre $S$.
Se obtiene:

\result{COROLARIO}{7.4} Sean $F,G$ esquemas en grupos de tipo finito sobre un
cuerpo $k$, y sea $u:F\to G$ un homomorfismo de $k$-grupos. Entonces $u$ se
factoriza como $F\to F'\to G$, donde $F\to F'$ es el homomorfismo cociente
por el subgrupo cerrado $\operatorname{Ker}u$ de $F$, y $F'\to G$ es un
homomorfismo de grupos que es una inmersión cerrada. Existe el cociente
$G/F=G/F'$. El formalismo usual (del tipo proporcionado por los teoremas de
Noether) es válido entre grupos algebraicos sobre $k$.

Este resultado permite tratar de manera uniforme tanto el paso al cociente
entre grupos algebraicos en el sentido clásico (es decir, los que son
irreducibles sobre $k$ y lisos sobre $k$) como el paso al cociente por
subgrupos «infinitesimales»
\src{116}
considerados por CARTIER. Resulta ventajoso considerar las «hiperálgebras»
introducidas por CARTIER, a raíz de los trabajos de DIEUDONNÉ sobre los grupos
formales, como grupos en la categoría de los esquemas formales sobre $k$ y,
cuando proceda (si corresponden a hiperálgebras de rango finito sobre $k$),
como grupos algebraicos finitos sobre $k$.

\fsection{8}{Una conjetura}

Responde a la necesidad de formar cocientes por el grupo proyectivo que actúa
en ciertos subesquemas de «esquemas de Hilbert» (que sustituyen a las
variedades de Chow en la teoría de esquemas).

Sea $S$ un preesquema y sea $n$ un entero. A todo preesquema $S'$ sobre $S$
se asocia el grupo $\operatorname{Gl}(n,\Gamma(S',\mathcal O_{S'}))$ de
matrices invertibles de tamaño $n$ por $n$ con entradas en el anillo de
secciones de $\mathcal O_{S'}$. Esto da un funtor contravariante en $S'$, que
se comprueba fácilmente que es representable; corresponde, por tanto, a un
esquema en grupos sobre $S$, de hecho afín sobre $S$, denotado por
$\operatorname{Gl}(n)_S$. Su formación es compatible con el cambio de base,
de modo que, en realidad, todo procede de un esquema en grupos sobre
$\mathbf Z$, denotado por $\operatorname{Gl}(n)$. El grupo
$\operatorname{Gl}(1)$, llamado grupo multiplicativo y denotado a menudo por
$\mathbf G_m$, corresponde al funtor
$S\longmapsto\Gamma(S,\mathcal O_S)^*$, el grupo de las «unidades» sobre $S$.
Hay un homomorfismo evidente
$\operatorname{Gl}(1)\to\operatorname{Gl}(n)$, y se construye fácilmente el
grupo cociente, denotado por $\operatorname{GP}(n-1)$ y llamado \emph{grupo
proyectivo de grado $n-1$ sobre $\mathbf Z$}. Representa el funtor que asocia
a $S$ el grupo de secciones del haz
$\widetilde{\operatorname{Gl}(n)}_S/
\widetilde{\operatorname{Gl}(1)}_S$, donde
$\widetilde{\operatorname{Gl}(n)}_S$ denota el haz de gérmenes de secciones de
$\operatorname{Gl}(n)_S$ sobre $S$. (¡Hay que tener cuidado: las secciones de
$\operatorname{GP}(n-1)_S$ sobre $S$ no proceden, en general, de secciones de
$\operatorname{Gl}(n)_S$ sobre $S$!) Obsérvese que puede demostrarse que
$\operatorname{GP}(n-1)$ representa también el funtor
$S\longmapsto\operatorname{Aut}_S(\mathbf P_S^{n-1})$, donde
$\mathbf P_S^{n-1}$ es el espacio proyectivo estándar de dimensión relativa
$n-1$ sobre $S$, al menos para $S$ noetheriano. Así es como interviene en la
teoría de módulos.

Sea $S$ un esquema noetheriano, que si se desea puede suponerse afín, y sea
$X$ un $S$-preesquema cuasiproyectivo provisto de un haz invertible
$\mathcal L$ muy amplio con respecto a $S$. Supongamos que el grupo
$G=\operatorname{GP}(n)_S$ actúa en $X$ y al mismo tiempo en $\mathcal L$
(compatiblemente con su acción en $X$), y que actúa \emph{libremente} en
$X$.\footnote{El texto impreso dice «en $S$». La acción considerada y la
condición de cociente conciernen a $X$.}
\src{117}

\result{CONJETURA}{8.1}\footnote{La fe de erratas vinculada de la Exposición 212
afirma que esta conjetura es definitivamente falsa. También señala que el
hecho positivo de la Observación 8.1 fue demostrado posteriormente y remite
al teorema de Mumford y a refinamientos posteriores.} En las condiciones
anteriores:

\smallskip
\noindent $1^\circ$ \emph{La relación de equivalencia definida por $G$ es
efectiva, el cociente $Y=X/G$ es de tipo finito sobre $S$, y el morfismo
canónico $f:X\to Y$ es plano y sobreyectivo} (así, $X$ se convierte en un
fibrado principal homogéneo sobre $Y$ con grupo
$G\times_S Y=\operatorname{GP}(n)_Y$).

\smallskip
\noindent $2^\circ$ Sea $\mathcal L'$ el haz invertible sobre $Y$ deducido de
$\mathcal L$ por «descenso fielmente plano» a lo largo de $f$ (véase [1], B,
\S~teo. 1). Entonces $\mathcal L'$ es \emph{preamplio sobre $Y$ con respecto
a $S$}, es decir, existen un entero $m$ y un morfismo cuasifinito de $Y$ a un
esquema proyectivo estándar adecuado $\mathbf P_S^N$ tales que
$(\mathcal L')^{\otimes m}$ es isomorfo a la imagen inversa de
$\mathcal O_{\mathbf P_S^N}(1)$.

Obsérvese que, incluso si $X$ es separado sobre $S$, puede ocurrir que $Y$ no
sea separado sobre $S$ (situación que se presenta en problemas de módulos
nada patológicos). Si se cumple $1^\circ$, entonces $Y$ es separado si y solo
si la relación de equivalencia definida por $G$ tiene grafo cerrado, es decir,
si $G\times X\to X\times X$ tiene imagen cerrada (es entonces una inmersión
cerrada). Si $Y$ es separado, entonces $\mathcal L'$ es preamplio sobre $Y$
con respecto a $S$ si y solo si es amplio, es decir, si una potencia tensorial
adecuada define una inmersión proyectiva. En los problemas de módulos
mencionados en la introducción puede demostrarse que la relación de
equivalencia obtenida tiene, en efecto, grafo cerrado.

\medskip
\noindent\textbf{OBSERVACIONES 8.1. ---} Hemos supuesto
$G=\operatorname{GP}(n)_S$ para fijar las ideas y porque este es actualmente
el caso más importante en la práctica. Una hipótesis razonable sobre $G$
parecería ser, en cambio, que $G$ fuese una de las «formas» sobre $S$ de uno
de los grupos de Tohokû (cuya construcción sobre los enteros fue realizada
también por CHEVALLEY). El único hecho positivo que conozco en la dirección de
la conjetura anterior es el siguiente: sea $X$ \emph{un esquema afín sobre un
cuerpo $k$ de característica $0$, en el que el grupo
$\operatorname{Gl}(n)_k$ o $\operatorname{GP}(n-1)_k$ actúa libremente.
Entonces la relación de equivalencia definida por $G$ es efectiva, el
cociente $X/G$ es también afín y el morfismo $X\to X/G$ es plano y
sobreyectivo}. La demostración utiliza el hecho siguiente (que actualmente
solo se ha demostrado en característica $0$): cuando se hace actuar a $G$ en
el anillo de coordenadas afín $A$ de $G$, considerado como espacio vectorial
sobre $k$, la representación trivial de $G$ aparece una sola vez (en una
serie de composición de un subespacio vectorial estable bajo $G$ y de
dimensión finita sobre $k$). Parece posible que un uso sistemático de
\src{118}
la teoría de representaciones lineales de $G$ conduzca finalmente a una
demostración de la conjetura, al menos sobre un cuerpo base. Cuando la base ya
no es un cuerpo, el conferenciante no sabe nada.

\section*{BIBLIOGRAFÍA}
\addcontentsline{toc}{section}{Bibliografía}

\begin{description}
\item[{[1]}] GROTHENDIECK (Alexander). --- Techniques de descente et théorèmes
d'existence en géométrie algébrique, I: Généralités, Séminaire Bourbaki,
1959/60, no. 190, 29 pp.

\item[{[2]}] GROTHENDIECK (Alexander). --- Techniques de descente et théorèmes
d'existence en géométrie algébrique, II: Le théorème d'existence en théorie
formelle des modules, Séminaire Bourbaki, 1959/60, no. 195, 22 pp.

\item[{[3]}] GROTHENDIECK (A.) and DIEUDONNÉ (J.). --- Éléments de géométrie
algébrique, I: Le langage des schémas. --- Paris, Presses universitaires de
France, 1960 (Institut des hautes Études scientifiques, Publications
mathématiques, 4).
\end{description}
